\chapter*{Arbitrary-Fs I2S dan IP komunikasi buatan sendiri}
\addcontentsline{toc}{chapter}{Arbitrary-Fs I2S dan IP komunikasi buatan sendiri}
\label{app:arbitrary-fs}

Lampiran ini menjelaskan varian desain I2S (\path{i2s_1_0}) yang menjadi
implementasi akhir penelitian. IP ini dibuat sebagai peripheral komunikasi buatan
sendiri agar sistem tidak bergantung pada IP berlisensi dari pihak ketiga.

\begin{itemize}
    \item Modul inti \path{i2s_slave_lite_v1_0_S00_AXI.v} membangkitkan \texttt{BCLK}, \texttt{WS}, dan \texttt{MCLK} menggunakan akumulator fase NCO pada domain \texttt{S\_AXI\_ACLK} 100 MHz.
    \item Frekuensi sampling diprogram melalui bitfield \texttt{SAMPLE\_RATE\_HZ} pada register \texttt{CONTROL} (rentang 0--1.048.575 Hz, dibatasi aman hingga 195.312 Hz).
    \item Register \texttt{STATUS[0]} berubah setiap awal frame stereo sehingga firmware dapat melakukan \textit{polling} tanpa interrupt.
    \item Port clock legacy \texttt{audio\_48\_clk}/\texttt{audio\_44\_clk} telah dihapus sepenuhnya dari modul level teratas (\path{i2s\_dds.v}) untuk menyederhanakan arsitektur.
    \item Kode lengkap berada di \path{ready_to_run/ip_repo/i2s_1_0/hdl}; firmware uji di \path{ready_to_run/sw/helloworld.c}.
\end{itemize}

Potongan kode NCO yang menunjukkan pembangkitan detak BCLK ditampilkan pada Listing~\ref{lst:app-dds-bclk}.

\begin{lstlisting}[language=Verilog, caption={Akumulator fase NCO untuk pembangkitan BCLK pada \texttt{i2s\_slave\_lite\_v1\_0\_S00\_AXI.v}.}, label={lst:app-dds-bclk}]
wire [31:0] bclk_toggle_rate_hz = {sample_rate_sync, 7'b0}; // Fs x 128
wire bclk_phase_wrap = (bclk_phase_sum >= AXI_CLK_HZ);

always @(posedge S_AXI_ACLK) begin
    if (bclk_phase_wrap) begin
        bclk_phase_acc_q <= bclk_phase_sum - AXI_CLK_HZ;
        bclk_q <= !bclk_q;
    end else begin
        bclk_phase_acc_q <= bclk_phase_sum;
    end
end
\end{lstlisting}

Perbedaan utama terhadap desain awal berbasis pembagi clock tetap:
\begin{itemize}
    \item frekuensi sampling tidak lagi dibatasi pada empat mode (\texttt{FS\_FAMILY}/\texttt{FS\_MODE}), melainkan dapat diprogram dalam Hz;
    \item serializer dan register file berada pada domain clock yang sama (100 MHz), sehingga tidak memerlukan CDC antar-domain untuk jalur data;
    \item \textit{handshake} firmware menggunakan register \texttt{STATUS}, bukan deadline berbasis \texttt{mcycle} atau \texttt{usleep()}.
\end{itemize}

Rincian kode dan perubahan perangkat lunak pendukung dijelaskan pada Bab IV dan Lampiran~\ref{app:code-summary}.
